% Conteggio requisiti:
% Funzionali Obbligatori: R1..R77-F-O = 77 requisiti
% Funzionali Desiderabili: R55-F-D = 1 requisito
% Funzionali Opzionali: R1..R6-F-Op = 6 requisiti
% Qualità Obbligatori: R1-Q-O, R2-Q-O, R3-Q-O = 3 requisiti
% Qualità Desiderabili: R1-Q-D = 1 requisito
% Vincolo Obbligatori: R1-V-O..R6-V-O = 6 requisiti
% Vincolo Opzionali: R1-V-Op = 1 requisito
% Prestazionali Obbligatori: R1-P-O = 1 requisito
% Prestazionali Desiderabili: R1-P-D = 1 requisito
% Totale Obbligatori: 77 + 3 + 6 + 1 = 87
% Totale Desiderabili: 1 + 1 + 1 = 3
% Totale Opzionali: 6 + 1 = 7

\section{Riepilogo}
Nella seguente sezione viene riportato un riepilogo quantitativo dello stato dei requisiti individuati, suddivisi in base al loro grado di importanza.

\subsection{Requisiti obbligatori soddisfatti}
\begin{table}[ht]
	\centering
	\renewcommand{\arraystretch}{1.4}
	\begin{tabularx}{0.8\textwidth}{|l|X|}
		\hline
		\rowcolor{primarycolor!10}
		\textbf{Tipologia Requisito} & \textbf{Numero Requisiti Obbligatori} \\
		\hline
		Funzionali & 77 \\
		\hline
		Di Qualità & 3 \\
		\hline
		Di Vincolo & 6 \\
		\hline
		Prestazionali & 1 \\
		\hline
		\textbf{Totale Obbligatori} & \textbf{87} \\
		\hline
	\end{tabularx}
	\caption{Riepilogo requisiti obbligatori}
\end{table}

\subsection{Requisiti desiderabili soddisfatti}
\begin{table}[ht]
	\centering
	\renewcommand{\arraystretch}{1.4}
	\begin{tabularx}{0.8\textwidth}{|l|X|}
		\hline
		\rowcolor{primarycolor!10}
		\textbf{Tipologia Requisito} & \textbf{Numero Requisiti Desiderabili} \\
		\hline
		Funzionali & 1 \\
		\hline
		Di Qualità & 1 \\
		\hline
		Di Vincolo & 0 \\
		\hline
		Prestazionali & 1 \\
		\hline
		\textbf{Totale Desiderabili} & \textbf{3} \\
		\hline
	\end{tabularx}
	\caption{Riepilogo requisiti desiderabili}
\end{table}

\subsection{Requisiti opzionali soddisfatti}
\begin{table}[ht]
	\centering
	\renewcommand{\arraystretch}{1.4}
	\begin{tabularx}{0.8\textwidth}{|l|X|}
		\hline
		\rowcolor{primarycolor!10}
		\textbf{Tipologia Requisito} & \textbf{Numero Requisiti Opzionali} \\
		\hline
		Funzionali & 6 \\
		\hline
		Di Qualità & 0 \\
		\hline
		Di Vincolo & 1 \\
		\hline
		Prestazionali & 0 \\
		\hline
		\textbf{Totale Opzionali} & \textbf{7} \\
		\hline
	\end{tabularx}
	\caption{Riepilogo requisiti opzionali}
\end{table}

\subsection{Riepilogo generale}
\begin{table}[ht]
	\centering
	\renewcommand{\arraystretch}{1.4}
	\begin{tabularx}{0.8\textwidth}{|l|X|}
		\hline
		\rowcolor{primarycolor!10}
		\textbf{Grado di Importanza} & \textbf{Totale Requisiti} \\
		\hline
		Obbligatori & 87 \\
		\hline
		Desiderabili & 3 \\
		\hline
		Opzionali & 7 \\
		\hline
		\textbf{Totale complessivo} & \textbf{97} \\
		\hline
	\end{tabularx}
	\caption{Riepilogo generale dei requisiti}
\end{table}